\section{Quantification}
\label{sec:quantification}

The calibration uses a stationary economy with four-year periods. A model
household represents an adult household unit. Housing quantities are measured
in occupied rooms, with housing services adjusted for tenure; financial
quantities are measured relative to model earnings.

\medskip
\noindent\fbox{\parbox{0.94\linewidth}{\textit{Provisional quantitative
contract.} The population counterpart, cohort alignment, within-period birth
timing, model-data observer, and working weights are under review.  In
particular, whether current earnings can be used at the purchase transaction is
an unresolved coordination item: the model text and an experimental calculation
use different timing conventions.  The earnings process and the utility
specification are also provisional inputs.  Thus the targets below define a
working calibration exercise rather than an adopted quantitative specification.}}

\subsection{Population, units, and external inputs}

The stationary economy aggregates household decisions over age, earnings, assets,
tenure, housing, and the number of children ever born and currently at home.
The model records completed fertility through children ever born and housing needs
through current dependents; these are distinct objects.  A four-year model period
therefore requires an explicit observer before an age- or event-time statistic can
be compared with its annual or biennial data counterpart.

The maintained earnings input combines an age--earnings profile with idiosyncratic
earnings risk.  Its parameters are estimated outside the household calibration,
after aggregating earnings to the model period.  This separation is useful because
earnings risk changes both resources at entry and the value of future housing and
fertility choices.  The measured earnings of the reference person and spouse
are represented by one exogenous household earnings process; household labor
supply is not a separate choice in this specification.

Entrants draw heterogeneous initial liquid wealth.  The working entry distribution
uses the nonhousing-wealth distribution of childless renters ages 18--24 as a
proxy for wealth at entry.  Its conversion from family-income ratios to the
model's earnings units and its joint distribution with earnings remain
provisional.  It is an external input to the
calibration, alongside the age profile, survival, fertility-success probabilities,
institutional tax and mortgage rules, and the housing menu.  These inputs determine
the feasible choices faced by a household before the parameters governing its
preferences and housing demand are selected.

\subsection{Calibration strategy}

Conditional on the external inputs, the calibration jointly selects the parameters
that govern intertemporal choice, the consumption--housing aggregator, the value of
owner-occupied services, the bequest motive, fertility costs and taste dispersion,
and the housing-supply intercept.  Alternative specifications may represent
children's housing needs through a minimum housing-service requirement or through
child-dependent housing expenditure shares.  The latter comparison is exploratory;
it is not a maintained utility choice.  Financing parameters and the age profile
are treated as institutional or external inputs.

Let $d_j(\vartheta)=\widehat m_j(\vartheta)-m_j$ denote the difference between
the model statistic and its empirical counterpart for moment $j$.  The working
criterion is
\begin{equation}
 \mathcal L(\vartheta)=\sum_{j\in\mathcal J}w_j d_j(\vartheta)^2,
 \label{eq:working-calibration-criterion}
\end{equation}
where $\mathcal J$ contains the scored rows in Table~\ref{tab:working-targets}.
The separate completed-fertility normalization is imposed and checked outside this
sum.  The $w_j$ are working discrepancy weights: they translate differences into a
common diagnostic scale, but do not yet constitute an efficient-SMM weighting
matrix or a statistical test of the model.  If $s_j$ denotes the working scale
reported in the table, then $w_j=1/s_j^2$.

The fertility block informs the level and timing of childbearing: it includes
childlessness, one-child status among mothers, the mean age at first birth, and the
share of first births at age 30 or older.  The saving and bequest block informs wealth
accumulation and transfers through aggregate wealth relative to gross earnings,
bequests relative to wealth, and old-age wealth dispersion.  The housing block
disciplines average occupied space, tenure, the change in rooms around a first
birth, the room difference across families with different numbers of resident
children, and the ownership difference for recent parents.  Taken together, these
blocks are intended to discipline the joint fertility, saving, tenure, and housing
choices rather than any one margin in isolation.

\subsection{Working moments and measurement}

Table~\ref{tab:working-targets} reports the complete working target set currently
under review.  Twelve rows are scored in the diagnostic criterion and one row is a
separate model normalization.  The reported scale identifies the unit used in the
working criterion; it is not reported for the normalization.  No model estimates
are reported in this table.

\begin{table}[htbp]
\centering
\caption{Working targets under review}
\label{tab:working-targets}
\small
\begin{tabularx}{\linewidth}{@{}X r r@{}}
\toprule
Moment & Target & Working scale \\
\midrule
Mean children ever born, ages 46+ (normalization) & 2.100000 & -- \\
Childless women, ages 40--44 & 0.198279 & 0.005305 \\
Exactly one child among mothers, ages 40--44 & 0.213655 & 0.006091 \\
Period mean age at first birth & 25.976264 & 0.084567 \\
First births at age 30+ & 0.249278 & 0.008492 \\
Wealth / annual gross labor earnings & 6.145861 & 0.362855 \\
Annual bequests / aggregate wealth & 0.008800 & 0.000440 \\
Old wealth/income p90 / median, ages 76--84 & 3.515935 & 0.306911 \\
Mean occupied rooms, capped at 9 & 5.561097 & 0.088381 \\
Ownership, heads ages 30--55 & 0.648334 & 0.020675 \\
First-birth room response, $-1$ to $+3$ & 0.720246 & 0.085260 \\
Rooms: 3+ versus 1--2 resident children & 0.347067 & 0.059705 \\
Recent-parent ownership gap & 0.162896 & 0.006080 \\
\bottomrule
\end{tabularx}
\begin{flushleft}\footnotesize
\textit{Notes:} Fertility stocks use the June 2004 and 2006 CPS supplements;
first-birth timing uses NCHS counts for 2003--2006. Wealth and old-age dispersion
use the 2003 and 2005 PSID waves. Housing levels and group differences use the
2005--2006 ACS in 42 metropolitan areas; the first-birth response uses a PSID
event study. The bequest restriction is external. Working scales combine
several conventions: approximate survey
sampling variation, temporal variation, bootstrap uncertainty, regression
uncertainty, and an externally assigned tolerance for bequests.  They are retained
for the present exercise while the target-and-observer contract is finalized.
\end{flushleft}
\end{table}

The fertility shares use pooled cross-sectional measures of children ever born,
whereas the timing statistics use first-birth flows.  The model counterpart must
therefore distinguish a stationary stock from a period flow and must state how a
birth within a four-year period is dated.  The completed-fertility normalization
is a model-side restriction on older households; it is not asserted to be the same
object as the CPS fertility cohorts in the table.

The wealth and bequest moments use explicit stock--flow denominators.  Aggregate
wealth is compared with annual gross labor earnings, and the old-age dispersion
uses a wealth-to-income ratio within its stated age range.  Initial wealth is
heterogeneous in the model, so these moments assess both lifecycle saving and the
distribution inherited at entry.  A common definition of wealth, income, and estate
coverage is required before these statistics can be interpreted as a fit.

Housing moments measure occupied rooms and owner occupancy over their stated
populations.  The first-birth room response is an event-time comparison in the
data, while the family-room and recent-parent moments are cross-sectional group
differences. Fit comparisons require consistent population and event-time rules,
complete moment and parameter tables, and lifecycle wealth, tenure and
housing-size profiles.
